\section{Conclusion}
\label{sec:discussion}

Our work puts forth a framework to examine the opinions reflected by LMs through the lens of public opinion polls. Using our \bname dataset, we identify a number of ways in which LMs are not well-aligned with human opinions, including overall representativeness with respect to people in the US; subgroup representativeness on groups such as 65+, Mormon, and widowed; and steerability. Our work also contributes to the broader discourse around LMs, including questions of whether instruct-tuning distorts opinion distributions, and whether models hold consistent biases towards liberal views.

\section{Limitations}

While our work provides a quantitative lens into LM opinions, it suffers from three key limitations discussed below.

\paragraph{Limitations of alignment.} Our approach analyzes LM opinions through the lens of \emph{who} they align with. This approach allows us to precisely define our metrics and collect data, but also warrants caution -- LMs that perfectly represent human opinions may not necessarily be desirable as they may also, in the process, replicate human biases. We view our metrics as useful ways to understand the behavior of LMs, and to allow model developers to identify when their models misrepresent specific groups, and not necessarily as benchmarks that should be blindly optimized.

\paragraph{Limitations of the ATP and surveys.} While our methodology to build an evaluation dataset can be easily adapted to any multiple-choice survey, the \bname dataset we instantiate is based on the American Trends Panel.
Surveys in general may be sensitive to details such as question specificity~\citep{berinsky2017measuring} and the ATP in particular has had past issues with social desirability bias~\citep{yan2021consequences}
that may affect the accuracy of the human opinion distribution. Beyond that, the ATP survey targets individuals in the US, making our conclusions valid only for the populations in the US. Many societies differ from WEIRD (Western, Educated, Industrialized, Rich and Democratic) societies such as the United States~\citep{henrich2010weirdest} and there is a need for future work on global equivalents to \bname.

\paragraph{Limitations of the multiple-choice format.} Our work focuses on probing LM behaviors using a multiple-choice prompt taken from public opinion surveys. While the multiple-choice format allows us to precisely evaluate the models, it also differs from the open-ended text generation setting in which LMs are being increasingly used. 
It is an open question whether opinion alignment that is measured through multiple choice will be reflected in the downstream use cases of LMs -- for example, will the liberal-leaning opinion alignment of RLHF fine-tuned models appear in a dialogue context or open-ended QA? Some recent works suggest that the group-alignment effects (e.g. to liberals) do reflect in other settings~\citep{perez2022discovering,hartmann2023political}, but whether these results transfer broadly warrants further investigation.
